\section{总结与延伸}

本讲把 Agent 从“技术热点”还原为“可经营系统”。讲者给出的核心启示是：企业级 Agent 的主战场不在单点模型能力，而在端到端可靠性、可验证结果、分层安全与持续优化机制。

\begin{table}[H]
\centering
\caption{全讲总结表：关键命题、证据与落地动作}
\begin{tabularx}{\textwidth}{>{\raggedright\arraybackslash}p{3.1cm} >{\raggedright\arraybackslash}X >{\raggedright\arraybackslash}p{4.2cm}}
\toprule
命题 & 讲者证据与观点 & 可执行落地动作 \\
\midrule
客服 Agent 的本质是结果系统 & 计费与评价围绕“是否解决问题”而非“是否聊得像人” & 以可验证动作定义 KPI，建立 outcome verification 事件模型 \\
单 Agent 将替代多渠道割裂 & 渠道正在从 phone/chat/email 向统一策略层收敛 & 统一状态机与身份/会话上下文，避免渠道 KPI 冲突 \\
Voice 是系统工程而非模型堆叠 & 延迟尾部、打断检测、噪声鲁棒决定体验成败 & 建立 turn-level latency budget、打断分类、provider failover \\
评测必须接近真实环境 & Towbench 强调工具交互、策略约束、多轮状态变化 & 以 simulation+回归集作为发布门禁，不仅依赖离线问答指标 \\
安全必须多层治理 & 提示注入与越权风险在真实流量中持续出现 & 输入/执行/输出三层 guardrail，加 deterministic access gateway \\
长期优势来自运营闭环 & tracing + insights + expert answers 构成持续学习系统 & 建立问题归因与知识回灌周会，形成季度迭代节奏 \\
\bottomrule
\end{tabularx}
\end{table}

\subsection{延伸阅读}

\begin{itemize}[leftmargin=1.5em]
    \item Towbench 项目与论文主页（关注 Tool-Agent-User 评测设定与 leaderboard）。
    \item Berkeley RDI Agentic AI MOOC F25 其余讲次（特别是 multi-agent、safety、infrastructure 相关课程）。
    \item OpenAI、Anthropic、Google 关于 agent evaluation、tool use、safety guardrails 的最新技术报告。
    \item 生产工程侧参考：service reliability、progressive delivery、observability 与 red-teaming 实践。
\end{itemize}

\subsection{给课程项目的实践建议}

如果你正在做课程项目，建议至少补齐以下四项：
第一，设计一个“可验证动作”而非仅文本回答的任务；
第二，构建最小多轮仿真评测集；
第三，加入输入与输出的双向守卫；
第四，记录完整 trace 并做一次系统性失败复盘。
做到这四项，你的 Agent 项目才真正跨过“演示系统”门槛，进入“可部署系统”范畴。

\end{document}
